% !TeX program = xelatex
\documentclass[11pt]{impeart}

% Das Handbuch verwendet bewusst einen globalen Hausstil. Gewöhnliche Dokumente
% sollten normalerweise das registrierte automatische Verhalten von \UseFont oder den Template-Schlüssel `fonts` verwenden.
\UseFont{libertinus}[global]
\setmonofont[
  BoldFont=lmmonolt10-bold.otf,
  ItalicFont=lmmono10-italic.otf,
  BoldItalicFont=lmmonolt10-boldoblique.otf
]{lmmono10-regular.otf}
\UseFunktionen{hyperlinks,tables}

\usepackage{catchfile}
\usepackage{metalogo}
\usepackage{pdfpages}
\usepackage[ngerman]{babel}

\newcommand{\IMPEVersionFile}{../../VERSION}
\newcommand{\IMPEShowcaseFile}{../showcase/impe-showcase.pdf}

\IfFileExists{\IMPEVersionFile}
  {\CatchFileEdef{\IMPEVersion}{\IMPEVersionFile}{\endlinechar=-1\relax}}
  {\PackageError{impe-manual}
     {Erforderliche Repository-Ressource ../../VERSION fehlt}
     {Aus manual/de kompilieren oder scripts/build_manual.ps1 aus dem Repository-Stammverzeichnis ausfuehren.}}

% Das erzeugte Handbuch bleibt reproduzierbar, wenn SOURCE_DATE_EPOCH festgelegt ist.
\ifdefined\XeTeXversion
  \special{pdf:trailerid [
    <5677bc84f2edee3de328d781062f390b>
    <778a1a55595b871c7ecbd03469f75deb>
  ]}
\else\ifdefined\pdftrailerid
  \expandafter\pdftrailerid\expandafter{IMPE-\IMPEVersion-manual-de}
\fi
\fi

\title{IMPE \LaTeX{} System}
\subtitle{Benutzerhandbuch}
\author{YANG Sikai}
\date{Version \IMPEVersion}

\newcommand{\IMPERepo}{https://github.com/KelvinYangBK67/IMPE-LaTeX-System}
\newcommand{\IMPEIssues}{https://github.com/KelvinYangBK67/IMPE-LaTeX-System/issues}
\newcommand{\IMPEShowcaseURL}{https://github.com/KelvinYangBK67/IMPE-LaTeX-System/blob/main/manual/showcase/impe-showcase.pdf}
\newcommand{\IMPETaggedShowcaseURL}{https://github.com/KelvinYangBK67/IMPE-LaTeX-System/blob/v\IMPEVersion/manual/showcase/impe-showcase.pdf}

\begin{document}

\maketitle

\begin{center}
\small Deutsche Übersetzung der englischen Referenzausgabe
\end{center}

\tableofcontents

\section{Einführung}

Das IMPE \LaTeX{} System ist ein modulares \XeLaTeX{}-Framework für
wiederverwendbare Layouts, Schriftrouting, mehrsprachigen Satz und optionale
Dokumentfunktionen.

Das System unterscheidet drei Bereiche, die sich unabhängig voneinander
konfigurieren lassen: Aus Sicht des Benutzers ist die wichtigste Unterscheidung die
zwischen \emph{Satzspiegeln}, \emph{Schriften} und \emph{Zusatzfunktionen}. Ein Layout
beschreibt die grundlegende visuelle Organisation eines Dokuments.
Schriftfamilien stellen globales oder lokales Schriftverhalten bereit,
einschließlich schriftsystembewussten Zuordnungs, wo dies erforderlich ist.
Funktionen stellen optionale Funktionen wie Mathematik, Hyperlinks, Zitate,
Tabellen oder Bilder bereit. Diese Teile können unabhängig voneinander
ausgewählt oder über eine einheitliche Einrichtungsschnittstelle kombiniert
werden.

Die normale Verwendung von IMPE ist bewusst knapp gehalten. Eine Wrapper-Klasse
kann ein sinnvolles Dokumentlayout und eine globale Standardschriftkonfiguration
festlegen; anschließend lädt der Benutzer die zusätzlichen Familien oder
Funktionen, die das Dokument benötigt. Wenn eine explizitere Kontrolle gewünscht
ist, lassen sich dieselben Komponenten in einem einzigen Aufruf von
\verb|\UseTemplateSet| deklarieren.

\XeLaTeX{} ist die primäre Dokument-Engine für IMPE. Diese Wahl ermöglicht es
dem System, sich auf Unicode-Text, OpenType-Schriften, \texttt{fontspec} und die
Inter-Character-Mechanismen von \XeTeX{} zu stützen und zugleich das vertraute
\LaTeX{}-Dokumentmodell beizubehalten. Für besondere Schriftsysteme ist zusätzlich die integrierte
XSR-Darstellung verfügbar. IMPE-Dokumente werden mit \XeLaTeX{} übersetzt.

Dieses Handbuch ist eine benutzerorientierte Einführung. Detaillierte
Dokumentation der Subsysteme wird im
Repository unter \texttt{docs/} gepflegt. Insbesondere die Dokumente zu
Schriften, Layouts, Funktionen und Systemarchitektur beschreiben
Registrierungsfelder und Implementierungsdetails, deren Wiederholung hier
unnötig wäre. Das englische Handbuch ist die Referenzausgabe; übersetzte
Handbücher sollen seiner Struktur und seinem technischen Inhalt folgen.

Eine visuelle Übersicht über das System bietet der Projekt-Showcase unter
\href{\IMPEShowcaseURL}{dem aktuellen IMPE-Showcase-PDF}. Der zugehörige Stand wird nach der Veröffentlichung unter dem
\href{\IMPETaggedShowcaseURL}{Versionstag v\IMPEVersion{}} bereitgestellt. Dieses kanonische
Repository-PDF wird außerdem in Anhang~\ref{app:showcase} wiedergegeben. Es soll
die Bandbreite an Schriftsystemen, Schreibrichtungen, Layouts und
Feature-Kombinationen demonstrieren, aus denen das System hervorgegangen ist.

\section{Installation}

\subsection{Ausgabeformen}

IMPE 1.1.0 wird als \texttt{Full}, \texttt{Core} und
\texttt{CTAN} vorbereitet. Alle drei Ausgaben stellen die
kanonischen \XeLaTeX{}-Klassen und dieselbe Kernfunktion bereit;
ihr Lieferumfang unterscheidet sich bei den Schriftdateien und
der Installation.

\begin{description}
  \item[\texttt{Full}] Enthält Laufzeit, Kompatibilitätseinträge,
    XSR und die ausdrücklich zur Weitergabe freigegebenen
    Schriftdateien des jeweiligen Pakets.
  \item[\texttt{Core}] Enthält Laufzeit, Kompatibilitätseinträge,
    XSR und Installationsprogramme. Schriften werden über
    \TeX{} Live, das Betriebssystem oder lokale Dateien aufgelöst.
  \item[\texttt{CTAN}] Das Archiv \texttt{impe-framework.zip} mit
    dem Verzeichnis \texttt{impe-framework/} umfasst die
    kanonischen Klassen, Laufzeit, XSR-Quellen und Unterlagen.
    Die Schriftdateien werden gesondert bereitgestellt.
\end{description}

Der Inhalt eines öffentlichen Full-Archivs richtet sich nach
einer geprüften Liste weitergabefähiger Schriftdateien.

\subsection{Full und Core installieren}

Das Installationsprogramm legt die verwaltete Laufzeit unter
\texttt{TEXMFHOME/tex/latex/impe/} ab. Als Einstiegspunkte
dienen \texttt{impe.sty}, \texttt{impeart.cls},
\texttt{impereport.cls}, \texttt{impebook.cls} und
\texttt{impebeamer.cls}. Die älteren Klassennamen stehen
in direkten Full/Core-Installationen weiterhin zur Verfügung.

\begin{verbatim}
# Windows, entpacktes Release:
.\install.bat
# Windows, Quellverzeichnis:
.\scripts\install.ps1

# Unix, entpacktes Release:
sh ./install.sh
# Unix, Quellverzeichnis:
sh ./scripts/install.sh
\end{verbatim}

Unter Windows akzeptiert das Installationsskript
\texttt{-TexmfRoot} und \texttt{-SkipXsrPython}; unter
Unix kann ein abweichender TEXMF-Pfad als Argument angegeben
werden. Falls die Aktualisierung der \TeX{}-Dateidatenbank
unterbleibt, kann \texttt{mktexlsr} nachträglich ausgeführt
werden. Bei einer Paketinstallation durch die jeweilige
\TeX{}-Distribution gelten deren eigene Verfahren.

\subsection{XSR und die Python-Laufzeit}

IMPE liefert XSR \textbf{1.3} mit. Die direkten
Installationsprogramme richten normalerweise eine eigene
virtuelle Python-Umgebung für die externe Darstellung ein.
Für dynamische ägyptische Hieroglyphen, die Stapelung
der kleinen Kitan-Schrift und Bildglyphen können
Python \textbf{3.11 oder neuer}, \texttt{venv} und
\texttt{pip} erforderlich sein. Die erstmalige Einrichtung
kann einen Internetzugang benötigen.

Eine reine \TeX{}-Installation überspringt die
Python-Einrichtung:

\begin{verbatim}
# Windows:
.\scripts\install.ps1 -SkipXsrPython
# Unix:
IMPE_SKIP_XSR_PYTHON=1 sh ./scripts/install.sh
\end{verbatim}

Für die dynamische XSR-Ausgabe benötigt \XeLaTeX{} einen
zulässigen Shell-Escape-Aufruf oder die vorgeschaltete
XSR-Verarbeitung. Außerdem müssen die tatsächlich
verwendeten TTF-/OTF-Dateien lesbar sein, damit der
Python-Zeichengenerator ihre Konturen laden kann.

\subsection{Schriftdateien und persönliche Einstellungen}

IMPE sucht eine registrierte Schrift zunächst als
lokale Datei und anschließend in den verfügbaren
\TeX{}- und Systemschriften. Auflösbare Systemschriften
müssen deshalb nicht in das Projektverzeichnis kopiert
werden. Für zusätzliche Schriftdateien dient
\texttt{assets/fonts/} als übliches Stammverzeichnis.

Eine dauerhafte Anpassung kann in
\texttt{impe.local.tex} stehen; auch eine direkte Angabe
im Vorspann ist möglich:

\begin{verbatim}
\SetCatalogFontRoot{D:/Fonts/IMPE}
\end{verbatim}

Bei nicht auflösbaren gewöhnlichen Schriften kann IMPE
nach Warnung eine Ersatzschrift verwenden. Die externe
XSR-Darstellung benötigt dagegen die wirkliche Schriftdatei.

\subsection{Ein kurzer Installationstest}

Auf einem Rechner mit Libertinus prüft dieses Beispiel
die Dokumentklasse, das Schriftverzeichnis und den lokalen
Schriftbefehl:

\begin{verbatim}
\documentclass{impeart}
\UseFont{libertinus}
\begin{document}
IMPE ist bereit. \LIB{Text mit Libertinus.}
\end{document}
\end{verbatim}

Der Test betrifft die übliche \XeLaTeX{}-Ausgabe;
für XSR und Veröffentlichung gibt es gesonderte Prüfungen.

\section{Schnellstart}

Der einfachste IMPE-Arbeitsablauf besteht darin, eine Wrapper-Klasse zu wählen,
die gewünschten Schriftfamilien oder Funktionen auszuwählen und das Dokument
anschließend wie gewohnt zu schreiben. IMPE ergänzt die gewöhnliche
\LaTeX{}-Dokumentstruktur um eine wiederverwendbare Einrichtungsschicht.

Ein kleiner Artikel kann daher beispielsweise so aussehen:

\begin{verbatim}
\documentclass[11pt]{impeart}

\UseFont{libertinus}
\UseFunktionen{math,hyperlinks}

\title{A Sample Document}
\subtitle{Typeset with IMPE}
\author{Author Name}

\begin{document}

\maketitle

\section{Introduction}
\LIB{Hello, world.}

\end{document}
\end{verbatim}

Die Klasse stellt bereits den englischen Satzspiegel und CMU bereit.
Der Aufruf \verb|\UseFont{libertinus}| registriert die lokale Auswahl
\verb|\LIB| und zugleich eine ausdrücklich gewählte globale
lateinische Schrift. Sie hat Vorrang vor der von der Klasse
vorgegebenen CMU-Schrift,
während \verb|\UseFunktionen| zwei optionale Funktionsmodule aktiviert. Der Rest
ist gewöhnliches \LaTeX{}.

Für ein Dokument, dessen Einrichtung sich besser als einzelne Deklaration
ausdrücken lässt, kann dieselbe Idee über \verb|\UseTemplateSet| geschrieben
werden:

\begin{verbatim}
\documentclass{impeart}

\UseTemplateSet{
  fonts = {libertinus},
  features = {math,hyperlinks}
}
\end{verbatim}

Ein Template-Set kann außerdem ein Layout und zusätzliche Schriftfamilien
auswählen. In der Praxis liefern IMPE-Klassen bereits ein geeignetes
Standardlayout, sodass die meisten Dokumente jene Einstellungen angeben,
die davon abweichen.

IMPE kann auch als Paket unterhalb einer standardmäßigen \LaTeX{}-Klasse
geladen werden:

\begin{verbatim}
\documentclass{article}
\usepackage{impe}

\UseTemplateSet{
  layout = en_doc,
  fonts = {libertinus},
  features = {hyperlinks}
}
\end{verbatim}

Diese Form ist nützlich, wenn die zugrunde liegende Klasse ausdrücklich
beibehalten werden muss oder wenn jede IMPE-Komponente manuell ausgewählt werden
soll. IMPE-Klassen und das rohe Paket verwenden dasselbe zugrunde liegende
System.

\section{Einsatzbereite Vorlagen}

Das Repository enthält unter \texttt{templates/} eine kleine Sammlung von
Starterprojekten. Sie sind für das tatsächliche Schreiben vorgesehen. Fixtures
für Regressionstests und interne Demonstrationen liegen unter \texttt{tests/}.

Die derzeitigen Startervorlagen sind:

\begin{description}
  \item[\texttt{article\_en/}] Englische Artikelvorlage mit \texttt{impeart}.
  \item[\texttt{article\_zh/}] Chinesische Artikelvorlage mit \texttt{impeart} und \texttt{lang=zh}.
  \item[\texttt{book\_en/}] Englische Buchvorlage mit \texttt{impebook}.
  \item[\texttt{book\_zh/}] Chinesische Buchvorlage mit \texttt{impebook} und \texttt{lang=zh}.
  \item[\texttt{beamer\_en/}] Englische Präsentationsvorlage mit \texttt{impebeamer}.
  \item[\texttt{beamer\_zh/}] Chinesische Präsentationsvorlage mit \texttt{impebeamer} und \texttt{lang=zh}.
\end{description}

Jede Vorlage folgt dem Stil eines installierten Pakets und enthält realistische
Metadaten, Gliederung und Dokumentinhalt sowie gegebenenfalls repräsentatives
Tabellen-, Abbildungs- oder Folienmaterial. Ein neues Projekt kann daher damit
beginnen, das nächstgelegene Vorlagenverzeichnis zu kopieren und den
Dokumentinhalt zu bearbeiten, anstatt einen Vorspann von Grund auf neu
aufzubauen.

Die Vorlagen verwenden bewusst die kanonischen \texttt{impe*}-Einstiegspunkte
und zeigen die normale benutzerseitige Praxis. Reine Entwicklungs-Fixtures
liegen unter \texttt{tests/}.

\section{Dokumentklassen und öffentliche Einstiegspunkte}

\subsection{Kanonische Klassen und Sprachwahl}

Für neue Dokumente stehen vier Klassen zur Verfügung:
\texttt{impeart}, \texttt{impereport}, \texttt{impebook}
und \texttt{impebeamer}. Die Option \texttt{lang=en}
(Standardeinstellung) oder \texttt{lang=zh} wählt die
Spracheinstellungen, ohne den Klassennamen zu ändern.

\begin{verbatim}
% Englischer Artikel
\documentclass[12pt,twoside]{impeart}

% Chinesischer Artikel
\documentclass[lang=zh,12pt]{impeart}

% Chinesisches Buch
\documentclass[lang=zh]{impebook}

% Chinesische Präsentation
\documentclass[lang=zh]{impebeamer}
\end{verbatim}

Den Klassen liegen \texttt{article}, \texttt{report},
\texttt{book} und \texttt{beamer} zugrunde.
Auf Englisch verwenden sie \texttt{en\_doc},
\texttt{en\_report}, \texttt{en\_book} und
\texttt{beamer}; auf Chinesisch \texttt{zh\_doc},
\texttt{zh\_report}, \texttt{zh\_book} und
\texttt{beamer\_zh}. CMU gehört zu den voreingestellten
Schriften; bei \texttt{lang=zh} kommt Shanggu hinzu.
Ausdrücklich vom Autor gewählte Schriften haben Vorrang.

Gewöhnliche Klassenoptionen wie \texttt{12pt},
\texttt{twoside} oder \texttt{aspectratio=169}
werden an die jeweilige Grundklasse weitergeleitet.
Die bisherigen \texttt{impe*\_zh}- und
\texttt{next*}-Klassennamen sind Kompatibilitätseinträge;
ihre Zuordnung steht in Abschnitt~\ref{sec:compatibility}.

\subsection{\texttt{impe.sty} direkt verwenden}

Die IMPE-Klassen teilen sich eine Implementierung und bieten komfortable
Standards. Ein Benutzer kann auch eine Standardklasse
beibehalten und

\begin{verbatim}
\usepackage{impe}
\end{verbatim}

laden, gefolgt von einer expliziten Konfiguration für Layout, Schriften und
Funktionen. Dies ist der empfohlene Weg, wenn die Dokumentklasse von einer
Zeitschrift, einem Verlag oder einem anderen Paket vorgegeben wird und als
Dokumentklasse bestehen bleiben muss.

\subsection{Titel und Untertitel}

IMPE unterstützt \verb|\subtitle{...}| neben den Standardbefehlen
\verb|\title| und \verb|\author|. IMPE-Klassen wenden ihre eigenen
Standardwerte für den Titelblock an, während der gewöhnliche
\LaTeX{}-Titelmechanismus erkennbar bleibt. Ein typischer Vorspann lautet
einfach:

\begin{verbatim}
\title{Main Title}
\subtitle{A More Specific Description}
\author{Author Name}
\end{verbatim}

gefolgt von \verb|\maketitle| im Dokumentkörper.

\section{Das Template-Modell}

\subsection{Vier Konfigurationsschlüssel}

Der wichtigste einheitliche Befehl ist \verb|\UseTemplateSet|. Sein
öffentliches Konfigurationsmodell ist bewusst klein gehalten:

\begin{verbatim}
\UseTemplateSet{
  layout = <preset>,
  fonts = {<families>},
  globalfonts = {<families>},
  features = {<features>}
}
\end{verbatim}

Die vier Schlüssel haben unterschiedliche Aufgaben.

\begin{description}
  \item[\texttt{layout}] wählt eine öffentliche Layout-Voreinstellung, die zur
    aktiven Dokumentklasse passt.
  \item[\texttt{fonts}] lädt Familien über ihr registriertes automatisches
    Verhalten. Dies ist die normale Wahl: Eine Familie kann einen lokalen
    Befehl definieren, eine dokumentweite Schrift anwenden oder
    bereichsspezifisches Zuordnung einrichten.
  \item[\texttt{globalfonts}] zwingt die aufgeführten Familien ausdrücklich
    durch ihren globalen Modus oder ihren globalen Zuordnung-Pfad.
  \item[\texttt{features}] lädt optionale Funktionsmodule.
\end{description}

\texttt{mainfonts} wird als Alias von \texttt{globalfonts} akzeptiert; beide
Schlüssel bieten explizite globale Überschreibungen; \texttt{fonts} bietet die
automatische Schnittstelle.

Die Einrichtungsschnittstelle bietet einen vorhersehbaren Ort für die vollständige
Konfiguration. IMPE-Klassen wenden bereits ihre eigenen Layout- und globalen
Schriftstandards an; ein Dokument kann daher allein eine zusätzliche
Feature-Liste angeben.

\subsection{Spezialisierte Kurzbefehle}

Für gezielte Konfiguration bietet IMPE öffentliche Kurzbefehle für die
einzelnen Subsysteme:

\begin{verbatim}
\UseLayout{en_doc}
\UseLayouts{...}

\UseFont{libertinus}
\UseFonts{arabic,tibetan}

\UseFont{libertinus}[global]
\UseLocalFonts{arabic,tibetan} % explicit local override
\UseGlobalFonts{libertinus}    % explicit global override

\UseFeature{hyperlinks}
\UseFunktionen{math,hyperlinks,tables}
\end{verbatim}

Die Font-Formen mit ausgelassenem Modus sind die empfohlenen Kurzbefehle für den normalen
Gebrauch. Die lokalen und globalen Aliase bleiben nützlich, wenn ein Dokument
das registrierte Verhalten einer Familie bewusst überschreibt. Die
vereinheitlichte und die spezialisierte Schnittstelle sind zwei Ansichten
desselben öffentlichen Modells.

\subsection{Komponierbare Templates}

Ein IMPE-Template kombiniert wiederverwendbare Komponenten. Ein Layout kann mit
unterschiedlichen Schriftauswahlen wiederverwendet werden, dieselbe
Feature-Menge kann in einem Artikel oder Buch zum Einsatz kommen, und eine
schriftspezifische Schriftfamilie kann ergänzt werden, während andere
Dokumentfunktionen erhalten bleiben. Diese Komposition ist das zentrale
Gestaltungsprinzip der öffentlichen Schnittstelle.

Die folgenden Deklarationen beschreiben beispielsweise zwei sehr
unterschiedliche Dokumente mit derselben Einrichtungsschnittstelle:

\begin{verbatim}
% Englischer Artikel: en_doc und CMU sind voreingestellt
\documentclass{impeart}
\UseTemplateSet{
  fonts = {libertinus},
  features = {math,hyperlinks}
}

% Chinesisches Buch: zh_book, CMU und Shanggu
\documentclass[lang=zh]{impebook}
\UseTemplateSet{
  fonts = {sanskrit,tibetan},
  features = {citations,index,hyperlinks}
}
\end{verbatim}

Die Beispiele sind getrennte Dokumentvorspänne. Die Klassen
wählen bereits Satzspiegel und Standardschriften; angegeben werden
nur weitere Schriften und Funktionen. Für abweichende Satzspiegel
steht der Schlüssel \texttt{layout=} zur Verfügung.

\section{Schriften, Schriftzuordnung und mehrsprachiger Satz}

Für die Schriftauswahl gibt es zwei öffentliche Wege: registrierte
Schriftvoreinstellungen für den alltäglichen Satz sowie ausdrücklich
zusammengesetzte Schreibprofile und Schriftressourcen. Beide
Schnittstellen sind in Version 1.1.0 verfügbar; die bekannten
lokalen Schriftbefehle bleiben erhalten.

\subsection{Registrierte Schriftvoreinstellungen}

\verb|\UseFont{id}| lädt eine registrierte Voreinstellung.
Auch der Schlüssel \texttt{fonts} in \verb|\UseTemplateSet|
nimmt solche Kennungen entgegen. Ohne zusätzliche Modusangabe
greift die automatische Registrierung der jeweiligen Familie:
Sie kann einen lokalen Befehl anlegen, eine globale Schrift
anmelden oder bestimmten Unicode-Bereichen eine Schrift zuweisen.

\begin{verbatim}
\UseFont{libertinus}
\UseFonts{sanskrit,tibetan,arabic}
\UseFont{libertinus}[local]
\UseFont{libertinus}[global]
\end{verbatim}

\verb|\UseFont{libertinus}| stellt \verb|\LIB| bereit und
meldet Libertinus zugleich als ausdrücklich gewählte
lateinische Hauptschrift an. Damit kann es die von einer
IMPE-Klasse vorgegebene CMU-Schrift ersetzen. Mit
\texttt{[local]} wird nur die lokale Bindung angelegt;
\texttt{[global]} verlangt eine global verwendbare
Voreinstellung. Für eine eindeutige Moduswahl gibt es
\verb|\UseLocalFont| und \verb|\UseGlobalFont|
sowie ihre Pluralformen; \verb|\UseMainFont| ist eine
weitere Bezeichnung für die globale Auswahl.

Lokale Befehle wie \verb|\SA{...}| übernehmen die
Schriftgestaltung innerhalb ihres Textarguments.
\verb|\Font{sanskrit}{...}| erlaubt die gleiche
ausdrückliche lokale Auswahl ohne vorausgehendes
\verb|\UseFont|. Die 59 registrierten Voreinstellungen
mit ihren lokalen Kurzbefehlen sind im Anhang verzeichnet.

\subsection{Schreibprofile und Schriftressourcen}

Das neue Schriftmodell unterscheidet vier Verzeichnisse:
Unicode-Blockdaten, \emph{Schreibprofile} für Zeichenbereiche
und Schriftsystemverhalten, \emph{Schriftressourcen}
für konkrete Schnitte sowie die registrierten
\emph{Voreinstellungen}. So bezeichnet
\texttt{devanagari} ein Schreibprofil, während
\texttt{sanskrit} eine bestehende Schriftvoreinstellung
mit besonderen Satzregeln ist.

IMPE 1.1.0 enthält 69 gepflegte Schreibprofile,
85 Ressourcenkennungen und 59 ursprüngliche
Voreinstellungen. Das Schreibprofil bestimmt die
Unicode-Abdeckung und die Darstellung; die Schriftressource
bestimmt die verwendeten Schriftdateien. Die Profile
\texttt{chinese}, \texttt{japanese} und \texttt{korean}
überschneiden sich bei Han-Zeichen; deren regionale
Glyphenform muss deshalb gegebenenfalls ausgewählt werden.

\begin{verbatim}
\UseWriting{[shanggu]chinese}
\UseWriting{[tiro]devanagari}
\UseWriting{[naskh]arabic}
\end{verbatim}

Ein Eintrag hat die Form
\verb|\UseWriting{[Ressource,Optionen]Profil}|.
Mehrere Einträge können außerhalb von eckigen Klammern
und geschweiften Gruppen durch Kommata getrennt werden.
Der erste Klammerwert kann eine eingetragene
Ressourcenkennung sein; \texttt{font=} wählt einen
auflösbaren Schriftnamen, \texttt{file=} eine konkrete
Schriftdatei. Pro Eintrag wird eine dieser Alternativen
gewählt. Weitere Optionen sind
\texttt{stylemap=native|full},
\texttt{writing=horizontal|vertical}, \texttt{offset=}
sowie einzelne Schriftstilzuweisungen wie
\texttt{italic=zhuque}. Eine globale Schriftauswahl
verwendet waagerechte Schreibrichtung; unterstützte
lokale Darstellungen können senkrecht gesetzt werden.

\subsection{Benannte und unmittelbare lokale Schriftauswahl}

\verb|\DeclareWriting{Name}{Konfiguration}| legt eine
wiederverwendbare Konfiguration an.
\verb|\Font{Name}[Optionen]{Text}| verwendet sie lokal;
\verb|\Writing{Konfiguration}{Text}| übernimmt die
unmittelbare Auswahl ohne Namen.

\begin{verbatim}
\DeclareWriting{han-note}
  {[shanggu,stylemap=full,italic=zhuque]chinese}
\begin{document}
\Font{han-note}{漢字}
\Font{sanskrit}{संस्कृतम्}
\Writing{[tiro]devanagari}{संस्कृतम्}
\end{document}
\end{verbatim}

Bei lokaler Auswahl gelten eine Schriftressource
und eine Darstellung für das gesamte Textargument.
Innerhalb davon kann eine weitere ausdrückliche lokale
Auswahl den äußeren Stil überschreiben. Auch
\texttt{font=} oder \texttt{file=} allein sind
als lokale Konfiguration möglich.

\texttt{stylemap=native} bevorzugt vorhandene
Schriftschnitte; \texttt{stylemap=full} erlaubt außerdem
ausdrücklich registrierte stilistische Ersetzungen.
Einige westliche Mehrschnittfamilien verwenden
standardmäßig die vollständige Zuordnung.
Eine fehlende Glyphe und ein fehlender Schriftschnitt
sind dabei unterschiedliche Fälle. Die einzelnen
Optionen stehen in \texttt{docs/USEWRITING-1.0.5.md};
die vollständigen Kennungen enthält der Anhang.

\subsection{Globale Priorität und Glyphenersatz}

Die Klassen stellen CMU und bei \texttt{lang=zh}
zusätzlich Shanggu als vorgegebene Schriften bereit.
Ausdrückliche Schriftauswahlen des Autors haben
Vorrang. Innerhalb derselben Prioritätsstufe bleibt
bei konkurrierenden lateinischen Hauptschriften
normalerweise die zuerst angemeldete maßgeblich.
Für überschneidende CJK-Bereiche gilt innerhalb
einer Stufe weiterhin Chinesisch vor Japanisch
vor Koreanisch.

\begin{verbatim}
\UseGlobalFonts{cmu,libertinus}
\end{verbatim}

Hier bleibt CMU die erste ausdrücklich gewählte
lateinische Hauptschrift; eine spätere Anmeldung von
Libertinus ersetzt sie nicht automatisch. Dagegen
kann \verb|\UseFont{libertinus}| die bloß implizite
CMU-Vorgabe der Klasse übersteuern.

Der automatische \emph{zeichenweise Glyphenersatz}
ist auf CJK beschränkt. Über xeCJK werden
Serifen/serifenlos $\times$ normal/fett
$\times$ aufrecht/kursiv (2$\times$2$\times$2)
nach verfügbaren Schnitten geprüft. Danach kommen
weitere registrierte Ressourcen desselben
Schreibprofils und schließlich andere CJK-Profile
in Betracht. Ausdrückliche Autorenauswahlen stehen
vor Klassenstandards, und neu angemeldete Ressourcen
aktualisieren die entsprechenden Ersatzketten.
Ein nicht auflösbarer Schriftname kann unabhängig
davon eine Ersatzschrift auslösen; für andere
Schriftsysteme bleibt die zusammenhängende
OpenType-Formung maßgeblich.

\subsection{Unicode-Bereiche und CJK-Schriften}

Globale Schriftsystembindungen können bestimmte
Unicode-Bereiche übernehmen; lokale Befehle
übersteuern diese Zuordnung in ihrem Textargument.
Für Han, Kana und Hangul nutzt IMPE die
CJK-Schriftkanäle von xeCJK. Gemeinsam verwendete
Han-Zeichen behalten die CJK-Hauptschrift;
die japanische oder koreanische Glyphenform lässt
sich beispielsweise durch \verb|\JP{...}| und
\verb|\KR{...}| gezielt wählen.

\subsection{Formung, Umbruch und Schreibrichtung}

Die OpenType-Formung komplexer Schriftsysteme
übernehmen \XeTeX{} und die Schriftdateien.
IMPE ergänzt passende Schrift-, Sprach- und
Umbrucheinstellungen, etwa für Sanskrit in
Devanagari, tibetische Tsheg-Grenzen und
thailändische Wörterbuchtrennung. Für entsprechende
Schriften stehen auch Rechts-nach-links-Funktionen
bereit. Mongolisch, Mandschurisch und Altuigurisch
besitzen lokale gedrehte senkrechte Satzformen.

\subsection{XSR 1.3: Hieroglyphen, Kitan und Bildglyphen}

XSR 1.3 ist in IMPE eingebunden. Es setzt ägyptische
Hieroglyphen, Stapel der kleinen Kitan-Schrift
sowie registrierte Bildglyphen. Die beiden Schriften
verwenden ihre eigenen Voreinstellungen;
Bildglyphen stehen über die Funktion \texttt{glyphs}
zur Verfügung.

\begin{description}
  \item[Ägyptisch] Nach \verb|\UseFont{egyptian}|
    können die Unicode-Steuerzeichen des Satzsystems
    oder \verb|\IMPEEgyptianText{...}| verwendet werden.
  \item[Kleine Kitan-Schrift] \verb|\UseFont{khitan_small}|
    aktiviert \verb|\KHS{...}| mit der nativen
    XSR-Stapelung. \verb|\KHSstack| und
    \verb|\KHSstackblock| stehen ebenfalls bereit.
    Clustertrenner und U+16FE4 bleiben erhalten;
    unterschiedliche Stapelhöhen werden ausgerichtet.
  \item[Bildglyphen] Mit \verb|\UseFeature{glyphs}|
    registriert \verb|\IMPEGlyphRegister| eine
    Bilddatei, die \verb|\IMPEGlyph| im Text setzt.
\end{description}

Für die dynamische Darstellung benötigt XSR Python,
Shell Escape oder Vorverarbeitung sowie lesbare
Schrift- und Bilddateien. Für hieroglyphische Quadrate
und Kitan-Cluster kann XSR den ursprünglichen linearen
Unicode-Text in PDF-\texttt{/ActualText} bewahren;
die tatsächliche Textübernahme hängt vom PDF-Leser ab.

\subsection{Weitere Hinweise zu Schriften}

Vertiefende Angaben stehen in \texttt{docs/FONTS.md}
und \texttt{docs/USEWRITING-1.0.5.md}.
Die Bildbeispiele enthält der Showcase.

\section{Satzspiegel und Zusatzfunktionen}

\subsection{Satzspiegel für unterschiedliche Veröffentlichungen}

Ein Satzspiegel legt Seitenmaße, Schriftgrößen,
Abstände und Kolumnentitel fest. IMPE bietet
für jede Dokumentart abgestimmte Einstellungen,
die sich mit gewöhnlichen \LaTeX{}-Paketen
weiter anpassen lassen.

\begin{description}
  \item[\texttt{en\_doc}, \texttt{zh\_doc}]
    Artikel auf A4 mit 29\,mm Seitenrand,
    28\,mm oben und 32\,mm unten.
  \item[\texttt{en\_report}, \texttt{zh\_report}]
    Eigenständige A4-Berichtsgestaltung mit
    31\,mm links, 29\,mm rechts, 29\,mm oben,
    34\,mm unten und eigenen Kolumnentiteln.
  \item[\texttt{en\_book}, \texttt{zh\_book}]
    Bücher im Format ISO B5 (176$\times$250\,mm),
    beidseitig, mit 19\,mm innerem und 31\,mm
    äußerem Rand sowie 3\,mm Bundstegzugabe.
  \item[\texttt{beamer}, \texttt{beamer\_zh}]
    Beamer mit dem Thema Madrid im Verhältnis 4:3.
\end{description}

Die Berichtsvorlagen haben eigene Seiteneinstellungen
und unterscheiden sich von den Artikelvorlagen.
Die wahlweise verfügbaren Satzspiegel
\texttt{beamer\_clean} und \texttt{beamer\_clean\_zh}
stellen die schlichte Präsentationsgestaltung
aus Version 1.0.6 bereit. Für Breitbildfolien kann
die Klassenoption \texttt{aspectratio=169} verwendet werden.

\subsection{Änderungen mit vorhandenen LaTeX-Befehlen}

Die kanonischen Artikel-, Berichts- und Buchklassen
laden \texttt{geometry} und \texttt{fancyhdr}.
Seitenränder und Kolumnentitel lassen sich
mit deren gewöhnlichen Befehlen nach der
Klassendeklaration anpassen:

\begin{verbatim}
\documentclass{impereport}
\geometry{a4paper,left=30mm,right=30mm}
\fancyhead[R]{Kurzer Kolumnentitel}
\end{verbatim}

Für einzelne registrierte Gestaltungswerte
dienen \verb|\LayoutSetup| und \verb|\LayoutSet|.
Mit \texttt{LayoutScope} können Absatzwerte
innerhalb des Dokuments vorübergehend geändert werden:

\begin{verbatim}
\documentclass{impeart}
\LayoutSetup{page.inner=25mm,page.top=24mm}
\begin{document}
\begin{LayoutScope}{text.par-indent=2em}
Ein Absatz mit lokalem Einzug.
\end{LayoutScope}
\end{document}
\end{verbatim}

Für die Seitengestaltung sind die regulären
\texttt{geometry}-Befehle, für Kolumnentitel
\texttt{fancyhdr} und für Folien die Beamer-Befehle
die bevorzugten Werkzeuge. Frühere Satzspiegel
lassen sich über die Kompatibilitätseinstellungen
in Abschnitt~\ref{sec:compatibility} auswählen.

\subsection{Verfügbare Zusatzfunktionen}

Zusatzfunktionen laden abgestimmte Pakete und
stellen deren reguläre \LaTeX{}-Befehle bereit:

\begin{description}
  \item[\texttt{math}] Mathematischer Satz, AMS-Pakete,
    Theoremumgebungen und Schriftwahl für Formeln.
  \item[\texttt{hyperlinks}] Verweise, PDF-Lesezeichen
    und stabile Sprungziele.
  \item[\texttt{citations}] Literaturverwaltung
    mit \texttt{biblatex} und \texttt{biber}.
  \item[\texttt{index}] Sachregister mit
    \texttt{imakeidx} und \texttt{xindy}.
  \item[\texttt{tables}] Tabellen mit
    \texttt{booktabs}, \texttt{longtable} und
    \texttt{tabularx}.
  \item[\texttt{image}] Abbildungen, Beschriftungen
    und Teilabbildungen.
  \item[\texttt{drawing}] Zeichnungen mit TikZ/PGF,
    \texttt{pgfplots} (Kompatibilitätsstufe 1.18)
    und \texttt{forest}.
  \item[\texttt{glyphs}] Registrierung und Satz
    einzelner Bildzeichen über XSR.
  \item[\texttt{headers}] Einstellbare Kolumnentitel
    mit \texttt{fancyhdr}.
\end{description}

\begin{verbatim}
\UseFunktionen{math,hyperlinks,tables,image,drawing,glyphs}
\end{verbatim}

Mehrere Funktionen können kombiniert werden.
Bereits aktivierte Funktionen werden nicht
ein zweites Mal geladen.

\subsection{Mathematik und Literaturangaben}

Die Funktion \texttt{math} stellt
\texttt{theorem}, \texttt{lemma},
\texttt{proposition}, \texttt{corollary},
\texttt{definition}, \texttt{example}
und \texttt{remark} bereit.
Die Formelschrift bleibt standardmäßig
Computer Modern; für Libertinus Math
kann vor dem Laden angegeben werden:

\begin{verbatim}
\UseMathFont{libertinus}
\UseFeature{math}
\end{verbatim}

Literaturangaben verwenden die gewohnten
\texttt{biblatex}-Befehle
\verb|\addbibresource|, \verb|\textcite|,
\verb|\parencite| und \verb|\printbibliography|.
Über \verb|\UseCitationStyle| lassen sich
APA, GB/T 7714, numerische oder
Autor-Jahr-Zitierweisen einstellen.

\subsection{Sachregister mit den ursprünglichen Befehlen}

Für neue Dokumente sind die regulären
\texttt{imakeidx}-Befehle \verb|\makeindex|,
\verb|\index| und \verb|\printindex|
vorgesehen. IMPE vergibt im selbst eingerichteten
Standardregister für jeden gewöhnlichen
\verb|\index|-Eintrag ein eigenes PDF-Sprungziel,
auch bei wiederholter Seitennummer oder mehreren
Einträgen auf einer Seite.

\begin{verbatim}
\UseFeature{index}
\begin{document}
Buddhismus\index{Buddhismus} und
Mahayana\index{Buddhismus!Mahayana}.
Später nochmals Buddhismus\index{Buddhismus}.
\printindex
\end{document}
\end{verbatim}

Hat der Benutzer bereits ein eigenes Standardregister
eingerichtet, bleibt dessen Konfiguration maßgeblich.
Weitere benannte Register können wie üblich
mit \verb|\makeindex[name=people,title={Personen},program=xindy]|
angelegt und durch \verb|\index[people]{...}|
gefüllt werden. Die individuellen Sprungziele
gelten für gewöhnliche Einträge im von IMPE
eingerichteten Standardregister. Besondere
Registersyntax wie \texttt{see}, \texttt{seealso},
Bereichsmarkierungen und benannte Register
behält ihre normalen Eigenschaften.

Der Aufruf für das von IMPE eingerichtete
Standardregister lautet im Ausgabeverzeichnis:

\begin{verbatim}
texindy -L english -C utf8 -M <job>-impe <job>.idx
\end{verbatim}

Danach wird \XeLaTeX{} zweimal ausgeführt.

\subsection{Tabellen, Abbildungen und Bildzeichen}

Neue Tabellen verwenden \texttt{table},
\texttt{tabular}, \texttt{longtable} und
die Befehle von \texttt{booktabs};
Abbildungen verwenden \texttt{figure},
\verb|\includegraphics| und
\texttt{subfigure}. Für Zeichnungen
stehen \texttt{tikzpicture},
\texttt{axis} und \texttt{forest}
bereit. Ältere Hilfsumgebungen für
Tabellen und Abbildungen bleiben in
bestehenden Dokumenten verwendbar.

\begin{verbatim}
\UseFeature{glyphs}
\IMPEGlyphRegister[logical={Zeichen},description={Abreibung}]
  {bs-042}{images/bs-042.svg}
\begin{document}
Vor \IMPEGlyph{bs-042} nach.
\IMPEGlyph[scale=1.2,raise=1pt,trim={1pt 0pt 1pt 0pt}]{bs-042}
\end{document}
\end{verbatim}

Die Bilddatei muss für die Darstellung verfügbar sein.
Die Metadaten \texttt{logical} und \texttt{description}
beschreiben das Zeichen; die Option \texttt{trim}
beschneidet den sichtbaren Ausschnitt.

\subsection{Kolumnentitel und Verknüpfungen}

\texttt{headers} kann einen verkürzten
Kolumnentitel aus der ersten Zeile von
\verb|\title| ableiten. Mit
\verb|\HeaderTitle{...}| wird ein
fester Titel angegeben; weitere
Einstellungen erfolgen über \texttt{fancyhdr}.

\texttt{hyperlinks} lädt \texttt{hyperref}
und \texttt{bookmark} mit Unicode-fähigen
PDF-Einstellungen. Die Befehle \verb|\label|,
\verb|\ref|, \verb|\href| und \verb|\url|
bleiben die reguläre Schnittstelle.
Weitere Einzelheiten stehen in
\texttt{docs/FEATURES.md}.

\section{Kompatibilität und Migration}\label{sec:compatibility}

Neue Dokumente verwenden die Klassen
\texttt{impeart}, \texttt{impereport},
\texttt{impebook}, \texttt{impebeamer}
oder das Paket \texttt{impe.sty};
für Chinesisch wird \texttt{lang=zh}
angegeben. Die früheren Klassen
\texttt{impeart\_zh}, \texttt{impereport\_zh},
\texttt{impebook\_zh} und \texttt{impebeamer\_zh}
entsprechen den jeweiligen kanonischen Klassen
mit \texttt{lang=zh}. Ebenso bleiben
\texttt{next*}-Klassen und \texttt{nextsystem.sty}
als Kompatibilitätseinträge verfügbar.

Die vorhandenen \verb|\UseFont|-Kennungen,
lokalen Befehle und ausdrücklich gewählten
Modi gelten weiterhin. Die neuen Schreibprofile
und Ressourcen sind eine zusätzliche,
eigenständige Auswahlmöglichkeit.

Für Sachregister werden die ursprünglichen
\verb|\index|-Befehle empfohlen.
\verb|\Term| und \verb|\IndexTitle|
bewahren ihre bisherigen Formatierungs-
und Erstvorkommnisregeln; gewöhnliche
\verb|\index|-Aufrufe können dagegen
jedes Vorkommen verzeichnen. Für eigentliche
Glossare kann das Paket \texttt{glossaries-extra}
direkt verwendet werden.

Auch die älteren Tabellen- und Abbildungshilfen
\texttt{TableInlineFit}, \texttt{TableLong},
\texttt{TableBook}, \texttt{TableBookX},
\texttt{TableBookNotes}, \texttt{NiceBooktable*},
\texttt{OneImage}, \texttt{OneImageInline},
\texttt{PanelFigure} und \verb|\Panel|
bleiben erhalten. Für neue Arbeiten sind
reguläre \LaTeX{}-Tabellen und Abbildungen
die erste Wahl; \texttt{ExampleBlock}
wird ebenfalls weiterhin unterstützt.

Mit \verb|\UseLegacyLayout| lassen sich
die Satzspiegelwerte von Version 1.0.5
für die jeweilige Klasse und Sprache
wiederherstellen:

\begin{verbatim}
\documentclass[lang=zh]{impebook}
\UseLegacyLayout
\end{verbatim}

Daneben bleiben die älteren
\texttt{report}- und \texttt{legacy\_*}-
Satzspiegelkennungen verfügbar.
Unterschiedliche Schriftversionen und
Paketstände können den Seitenumbruch
weiterhin beeinflussen.

Das Installationsprogramm erkennt
verwaltete Installationen der Reihe 0.1.x
unter \texttt{tex/latex/nextsystem/}.
Es legt die aktuelle Laufzeit unter
\texttt{tex/latex/impe/} ab,
übernimmt bekannte lokale Einstellungen
und entfernt ausschließlich erkannte
Altdateien der verwalteten Installation.
Für neue Einstellungen dient
\texttt{impe.local.tex}.

\section{Paketstruktur und Erweiterungspunkte}

Das Repository trennt öffentliche Einstiegspunkte, stabile Mechanismen,
Registrierungen, spezialisierte Implementierungen und lokale Ressourcen:

\begin{verbatim}
package/   öffentliche Paket- und Klasseneinstiege
core/      gemeinsame Mechanismen
catalog/   Kennungen und Registrierungsdaten
modules/   besondere Implementierungen
assets/    lokale Ressourcen wie Schriften
scripts/   Installation und Veröffentlichung
manual/    Handbuchquellen
docs/      ausführliche Systemdokumentation
\end{verbatim}

\subsection{Der stabile Kern}

Die Schicht \texttt{core/} enthält stabile gemeinsame Mechanismen für neue
Schriftfamilien, Layout-Voreinstellungen und Funktionen. Sie implementiert die
Schriftregistrierung und
Zuordnung-Mechanik, die Layout-Registrierung und Kompatibilitätsprüfungen, den
Funktionslader und damit verbundenes gemeinsames Verhalten.

Gewöhnliche OpenType-Schriften lassen sich über den Katalog
anmelden und zusätzliche Funktionen mit dem zentralen
Lademechanismus verbinden. Dadurch bleibt der Kern
von den konkreten Schrift- und Funktionsdaten getrennt.

\subsection{Kataloge}

Die Schicht \texttt{catalog/} vergibt öffentliche IDs und Metadaten. Sie enthält
die zentralen Registrierungen für Schriften, Layouts und Funktionen. Für viele
Erweiterungen ist dies der richtige Ausgangspunkt.

Eine gewöhnliche Schriftfamilie mit standardmäßigem
\texttt{fontspec}-Schriftformung benötigt, kann in der Regel vollständig als
Katalogregistrierung ausgedrückt werden. Ebenso wird eine öffentliche
Layout-Voreinstellung aus vorhandenen internen Komponenten zusammengesetzt,
und eine Funktionskennung verweist auf ihr Implementierungsmodul. Die zentrale
Ablage dieser Deklarationen macht die öffentliche Oberfläche an einem Ort
überprüfbar.

\subsection{Module}

Die Schicht \texttt{modules/} ist Funktionalität vorbehalten, die tatsächlich
eine eigene Implementierung benötigt: Funktionsmodule, spezialisierte
Schriftsystemunterstützung oder wiederverwendbare interne Layout-Komponenten.
Ein neues Modul sollte daher eine echte Verhaltenserweiterung darstellen;
Schriftnamen gehören in Katalogregistrierungen.

Diese Regel ist für mehrsprachige Unterstützung besonders wichtig. Die meisten
Schriftsysteme sollten gewöhnliches OpenType-Schriftformung und Katalogmetadaten
verwenden. Spezialisierte Module dienen Schriftsystemen oder
Rendering-Strategien mit Verhalten über den generischen Kern hinaus.

\subsection{IMPE erweitern}

Entwickler, die das System erweitern möchten, beginnen mit den detaillierten
Registrierungsschnittstellen in \texttt{docs/SYSTEM.md} und konsultieren
anschließend das relevante Subsystemdokument: \texttt{docs/FONTS.md},
\texttt{docs/LAYOUTS.md} oder \texttt{docs/FEATURES.md}.

Die allgemeine Regel ist einfach: Eine Registrierung ist einem neuen Mechanismus
vorzuziehen, und bei tatsächlich außergewöhnlichem Verhalten ist ein
spezialisiertes Modul einer Änderung des stabilen Kerns vorzuziehen.

\section{Weiterführende Dokumentation und Rückmeldungen}

Dieses Handbuch ist der normale Ausgangspunkt für die Verwendung von IMPE. Das
Repository enthält detailliertere technische Dokumentation für Benutzer, die
eine vollständige Referenz benötigen oder das System erweitern möchten.

\begin{description}
  \item[\texttt{docs/SYSTEM.md}] Architektur, öffentliche Einstiegsschichten,
    Template-Konfiguration, Release-Modell und Erweiterungsgrenzen.
  \item[\texttt{docs/FONTS.md}] Familienregistrierung, lokale und globale Modi,
    Zuordnung, Ersatz, Schriftformung, Schreibverhalten, vertikale Unterstützung und
    externe Darstellung.
  \item[\texttt{docs/LAYOUTS.md}] öffentliche Layout-Voreinstellungen,
    Kompatibilitätsziele und das Komponentenmodell, aus dem Voreinstellungen
    aufgebaut werden.
  \item[\texttt{docs/FEATURES.md}] öffentliche Funktionskennungen sowie Befehle,
    Umgebungen, Paket-Stacks und Kompilierungsanforderungen der einzelnen
    Module.
  \item[\texttt{CHANGELOG.md}] Änderungen in veröffentlichten Versionen.
  \item[Showcase] der vollständige visuelle Showcase ist im
    \href{\IMPEShowcaseURL}{Projekt-Repository} verfügbar.
\end{description}

Das Projekt-Repository ist erreichbar unter

\begin{center}
\href{\IMPERepo}{\texttt{github.com/KelvinYangBK67/IMPE-LaTeX-System}}
\end{center}

Fehlerberichte, Funktionswünsche, Dokumentationskorrekturen,
Portabilitätsberichte und andere konkrete Rückmeldungen sind über GitHub Issues
willkommen:

\begin{center}
\href{\IMPEIssues}{\texttt{github.com/KelvinYangBK67/IMPE-LaTeX-System/issues}}
\end{center}

Bei der Meldung eines Satzproblems ist es hilfreich, eine minimale Quelldatei,
die \TeX{}-Engine und \TeX{}-Live-Version, die IMPE-Version und die relevanten
Schriftinformationen anzugeben. Berichte zu einer lokalen oder privat
bereitgestellten Schrift sollten, soweit möglich, außerdem zwischen einem
Problem der Glyphenabdeckung und einem IMPE-Zuordnung-Problem unterscheiden.

Das englische Handbuch dient als fachliche Referenz. Die deutsche
Fassung verwendet die in der deutschsprachigen Satztechnik üblichen
Bezeichnungen; Befehle und ihre Funktionsweise entsprechen der Referenzausgabe.

\appendix

\section{Kurzreferenz}\label{app:quick-reference}

Dieser Anhang enthält die in Version 1.1.0 verfügbaren
öffentlichen Schnittstellen und Kennungen. Zunächst
folgen die empfohlenen Klassen und Befehle,
danach die Schriftverzeichnisse und zuletzt
die Kompatibilitätseinträge.

\subsection{Kanonische Dokumentklassen}

Die vier Klassen verwenden \texttt{lang=en}
(Standardeinstellung) oder \texttt{lang=zh}.
Eine gesonderte Klasse für Chinesisch ist für
neue Dokumente nicht erforderlich.

\begin{longtable}{@{}p{0.25\linewidth}p{0.33\linewidth}p{0.34\linewidth}@{}}
\toprule
\textbf{Klasse} & \textbf{Englisch} & \textbf{Mit \texttt{lang=zh}} \\
\midrule
\endfirsthead
\toprule
\textbf{Klasse} & \textbf{Englisch} & \textbf{Mit \texttt{lang=zh}} \\
\midrule
\endhead
\texttt{impeart} & \texttt{en\_doc}; A4; CMU & \texttt{zh\_doc}; A4; CMU und Shanggu \\
\texttt{impereport} & \texttt{en\_report}; A4; CMU & \texttt{zh\_report}; A4; CMU und Shanggu \\
\texttt{impebook} & \texttt{en\_book}; B5; CMU & \texttt{zh\_book}; B5; CMU und Shanggu \\
\texttt{impebeamer} & \texttt{beamer}; Madrid 4:3; CMU & \texttt{beamer\_zh}; Madrid 4:3; CMU und Shanggu \\
\bottomrule
\end{longtable}

Bei einer durch einen Verlag vorgegebenen
Dokumentklasse kann stattdessen \verb|\usepackage{impe}|
geladen werden.

\subsection{Einrichtung und Schriftvoreinstellungen}

\begin{longtable}{@{}p{0.43\linewidth}p{0.49\linewidth}@{}}
\toprule
\textbf{Befehl} & \textbf{Zweck} \\
\midrule
\endfirsthead
\toprule
\textbf{Befehl} & \textbf{Zweck} \\
\midrule
\endhead
\texttt{\string\UseTemplateSet} & Satzspiegel, Schriften und Funktionen zusammen auswählen. \\
\texttt{\string\UseFont}, \texttt{\string\UseFonts} & Voreinstellungen automatisch oder mit \texttt{[local]} und \texttt{[global]} laden. \\
\texttt{\string\UseLocalFont}, \texttt{\string\UseLocalFonts} & Ausschließlich lokale Schriftauswahl. \\
\texttt{\string\UseGlobalFont}, \texttt{\string\UseGlobalFonts} & Ausdrücklich globale Schriftauswahl. \\
\texttt{\string\UseMainFont}, \texttt{\string\UseMainFonts} & Synonyme der globalen Auswahl. \\
\texttt{\string\UseLayout}, \texttt{\string\UseLayouts} & Eingetragene Satzspiegel wählen. \\
\texttt{\string\UseFeature}, \texttt{\string\UseFunktionen} & Zusatzfunktionen aktivieren. \\
\texttt{\string\SetCatalogFontRoot} & Verzeichnis lokaler Schriftdateien festlegen. \\
\texttt{\string\subtitle} & Untertitel des Dokuments setzen. \\
\bottomrule
\end{longtable}

Die Schlüssel von \verb|\UseTemplateSet| heißen
\texttt{layout}, \texttt{fonts}, \texttt{globalfonts}
und \texttt{features}; \texttt{mainfonts} bezeichnet
ebenfalls die globale Schriftauswahl.
Ausdrückliche Benutzerwahlen stehen über
Klassenstandards, innerhalb derselben Stufe
ist bei lateinischen Hauptschriften die
Anmeldereihenfolge maßgeblich.

\subsection{Schreibprofile und lokale Texte}

\begin{verbatim}
\UseWriting{[shanggu]chinese}
\DeclareWriting{han-note}{[shanggu,italic=zhuque]chinese}
\Font{han-note}{漢字}
\Writing{[tiro]devanagari}{संस्कृतम्}
\end{verbatim}

\begin{longtable}{@{}p{0.44\linewidth}p{0.48\linewidth}@{}}
\toprule
\textbf{Befehl} & \textbf{Zweck} \\
\midrule
\endfirsthead
\toprule
\textbf{Befehl} & \textbf{Zweck} \\
\midrule
\endhead
\texttt{\string\UseWriting} & Unicode-Bereiche einem Schreibprofil und einer Schriftressource zuweisen. \\
\texttt{\string\DeclareWriting} & Benannte Konfiguration speichern. \\
\texttt{\string\Font} & Benannte Voreinstellung oder Konfiguration lokal verwenden. \\
\texttt{\string\Writing} & Unbenannte lokale Konfiguration anwenden. \\
\bottomrule
\end{longtable}

Die Konfiguration verwendet
\texttt{[Ressource,Optionen]Profil}; die
Optionsnamen bleiben unverändert:
\texttt{font=}, \texttt{file=},
\texttt{stylemap=native|full}, \texttt{writing=}
und die einzelnen Schriftstilzuweisungen.
Die vollständigen Profile und Ressourcen
stehen in den Tabellen weiter unten.

\subsection{Satzspiegel und Kolumnentitel}

\begin{longtable}{@{}p{0.36\linewidth}p{0.56\linewidth}@{}}
\toprule
\textbf{Kennung oder Befehl} & \textbf{Zweck} \\
\midrule
\endfirsthead
\toprule
\textbf{Kennung oder Befehl} & \textbf{Zweck} \\
\midrule
\endhead
\texttt{en\_doc}, \texttt{zh\_doc} & A4-Artikelgestaltung. \\
\texttt{en\_report}, \texttt{zh\_report} & Eigenständige A4-Berichtsgestaltung. \\
\texttt{en\_book}, \texttt{zh\_book} & Bücher im ISO-B5-Format. \\
\texttt{beamer}, \texttt{beamer\_zh} & Madrid, Folienformat 4:3. \\
\texttt{beamer\_clean}, \texttt{beamer\_clean\_zh} & Wahlweise schlichte Foliengestaltung. \\
\texttt{\string\LayoutSetup}, \texttt{\string\LayoutSet} & Registrierte Satzspiegelwerte festlegen. \\
\texttt{LayoutScope} & Absatzbezogene Einstellungen. \\
\texttt{\string\HeaderStyle}, \texttt{\string\HeaderTitle} & Form und festen Text von Kolumnentiteln einstellen. \\
\bottomrule
\end{longtable}

Für einzelne Ränder und Kolumnentitel werden
gewöhnlich \verb|\geometry| und \texttt{fancyhdr}
verwendet.

\subsection{Kennungen der Zusatzfunktionen}

\begin{longtable}{@{}p{0.25\linewidth}p{0.68\linewidth}@{}}
\toprule
\textbf{Kennung} & \textbf{Reguläre Schnittstelle} \\
\midrule
\endfirsthead
\toprule
\textbf{Kennung} & \textbf{Reguläre Schnittstelle} \\
\midrule
\endhead
\texttt{math} & AMS-Mathematik, Theoremumgebungen, \texttt{\string\UseMathFont}. \\
\texttt{hyperlinks} & \texttt{hyperref}, \texttt{bookmark}, Verweise und Lesezeichen. \\
\texttt{citations} & \texttt{biblatex}, Literaturverzeichnis. \\
\texttt{index} & \texttt{\string\makeindex}, \texttt{\string\index}, \texttt{\string\printindex}. \\
\texttt{tables} & \texttt{tabular}, \texttt{longtable}, \texttt{booktabs}. \\
\texttt{image} & \texttt{\string\includegraphics}, Teilabbildungen. \\
\texttt{drawing} & TikZ, \texttt{axis}, \texttt{forest}. \\
\texttt{glyphs} & \texttt{\string\IMPEGlyphRegister}, \texttt{\string\IMPEGlyph}. \\
\texttt{headers} & Kolumnentitel über \texttt{fancyhdr}. \\
\bottomrule
\end{longtable}

Schriften für ägyptische Hieroglyphen und die
kleine Kitan-Schrift aktiviert man über
\verb|\UseFont{egyptian}| beziehungsweise
\verb|\UseFont{khitan_small}|.

\subsection{Registrierte Schriftvoreinstellungen (59)}

Die folgenden 59 Familienkennungen werden von
\verb|\UseFont| erkannt. Der Eintrag unter
\textbf{Auto} beschreibt die Wirkung
bei ausgelassener Modusangabe:
\texttt{G} steht für globale Anmeldung,
\texttt{L} für lokale Befehle,
\texttt{L+G} für beides und
\texttt{L+R} für lokale Befehle mit
Unicode-Bereichszuordnung. Die tatsächliche
Hauptschrift hängt außerdem von Priorität
und Anmeldereihenfolge ab.

\begingroup\small
\begin{longtable}{@{}p{0.285\linewidth}p{0.10\linewidth}p{0.17\linewidth}p{0.33\linewidth}@{}}
\toprule
\textbf{Familienkennung} & \textbf{Auto} & \textbf{Lokaler Befehl} & \textbf{Verwendung} \\
\midrule
\endfirsthead
\toprule
\textbf{Familienkennung} & \textbf{Auto} & \textbf{Lokaler Befehl} & \textbf{Verwendung} \\
\midrule
\endhead
\texttt{cmu} & G & --- & Computer Modern Unicode als Grundschrift \\
\texttt{noto} & L+G & \texttt{\string\NOT} & Lateinische Noto-Schriftfamilie \\
\texttt{times} & L+G & \texttt{\string\TIM} & Times mit Arial und Consolas \\
\texttt{gentium} & L+G & \texttt{\string\GEN} & Gentium Plus \\
\texttt{charis} & L+G & \texttt{\string\CHA} & Charis SIL \\
\texttt{libertinus} & L+G & \texttt{\string\LIB} & Libertinus Serif, Sans und Mono \\
\texttt{mlmodern} & L & \texttt{\string\MLM} & Lokale Schrift MLModern \\
\texttt{anatolian} & L & \texttt{\string\CA} & Karische Schrift \\
\texttt{coptic} & L & \texttt{\string\CO} & Koptische Schrift \\
\texttt{bopomofo} & L & \texttt{\string\ZY} & Zhuyin-Lautzeichen \\
\texttt{cuneiform} & L & \texttt{\string\CU} & Keilschrift \\
\texttt{glagolitic} & L & \texttt{\string\GL} & Glagolitische Schrift \\
\texttt{italic} & L & \texttt{\string\OI} & Altitalische Schrift \\
\texttt{hungarian} & L & \texttt{\string\OH} & Altungarische Schrift \\
\texttt{runic} & L & \texttt{\string\RU} & Runenschrift \\
\texttt{armenian} & L & \texttt{\string\HY} & Armenische Schrift \\
\texttt{georgian} & L & \texttt{\string\KA} & Georgische Schrift \\
\texttt{hindi} & L+R & \texttt{\string\HI} & Devanagari mit Hindi-Zeilenumbruch \\
\texttt{sanskrit} & L+R & \texttt{\string\SA} & Devanagari mit Sanskrit-Zeilenumbruch \\
\texttt{devanagari} & L & \texttt{\string\DEV} & Allgemeines Devanagari \\
\texttt{tamil} & L & \texttt{\string\TA} & Tamilische Schrift \\
\texttt{brahmi} & L & \texttt{\string\BR} & Brahmi-Schrift \\
\texttt{tibetan} & L+R & \texttt{\string\TI} & Tibetisch mit Tsheg-Umbruch \\
\texttt{segoe} & L & \texttt{\string\SEG} & Segoe UI Historic \\
\texttt{thai} & L & \texttt{\string\TH} & Thailändische Schrift \\
\texttt{arabic} & L & \texttt{\string\AR} & Arabisch mit verschiedenen Schriftschnitten \\
\texttt{urdu} & L & \texttt{\string\UR} & Urdu in Nastaliq \\
\texttt{aramaic} & L & \texttt{\string\IA} & Reichsaramäische Schrift \\
\texttt{nabataean} & L & \texttt{\string\NB} & Nabatäische Schrift \\
\texttt{hebrew} & L & \texttt{\string\HE} & Hebräische Schrift \\
\texttt{syriac} & L & \texttt{\string\SY} & Syrische Schrift \\
\texttt{syriac\_eastern} & L+R & \texttt{\string\SYE} & Ostsyrische Schrift mit Bereichszuordnung \\
\texttt{kharosthi} & L & \texttt{\string\KH} & Kharoshthi-Schrift \\
\texttt{egyptian} & L & \texttt{\string\EG} & Ägyptische Hieroglyphen über XSR \\
\texttt{khitan\_small} & L & \texttt{\string\KHS} & Kleine Kitan-Schrift mit XSR-Stapelung \\
\texttt{pahlavi\_parthian} & L+R & \texttt{\string\PAR} & Inschriftenparthisch \\
\texttt{pahlavi\_inscriptional} & L+R & \texttt{\string\PAH} & Inschriftenpahlavi \\
\texttt{pahlavi\_psalter} & L & \texttt{\string\PSP} & Psalterpahlavi \\
\texttt{avestan} & L & \texttt{\string\AV} & Avestische Schrift \\
\texttt{manichaean} & L & \texttt{\string\MA} & Manichäische Schrift \\
\texttt{phoenician} & L & \texttt{\string\PH} & Phönizische Schrift \\
\texttt{samaritan} & L & \texttt{\string\SM} & Samaritanische Schrift \\
\texttt{sogdian} & L & \texttt{\string\SG} & Sogdische Schrift \\
\texttt{sogdian\_old} & L+R & \texttt{\string\SGO} & Altsogdische Schrift mit Bereichszuordnung \\
\texttt{chinese\_simplified} & L+R & \texttt{\string\SC} & Vereinfachte chinesische Schrift \\
\texttt{chinese\_traditional} & L+R & \texttt{\string\TC} & Traditionelle chinesische Schrift \\
\texttt{japanese} & L+R & \texttt{\string\JP} & Japanische Schrift und Kana \\
\texttt{wenjin} & L+R & \texttt{\string\WJ} & WenJin Mincho mit CJK-Glyphenersatz \\
\texttt{shanggu} & G & --- & Globale CJK-Schrift Shanggu \\
\texttt{sim} & G & --- & Globale Systemschrift für vereinfachtes Chinesisch \\
\texttt{korean} & L+R & \texttt{\string\KR} & Koreanische Schrift und Hangul \\
\texttt{tangut} & L+R & \texttt{\string\TG} & Tangutische Schrift mit CJK-Zuordnung \\
\texttt{mongolian} & L & \texttt{\string\MO} & Mongolisch mit lokaler Senkrechtschrift \\
\texttt{mongolian\_baiti} & L & \texttt{\string\MOb} & Mongolisch Baiti, senkrecht \\
\texttt{manchu} & L & \texttt{\string\MC} & Mandschurisch, senkrecht \\
\texttt{turkic} & L & \texttt{\string\OT} & Alttürkische Schrift \\
\texttt{uyghur} & L & \texttt{\string\UY} & Altuigurische Schrift mit Drehung \\
\texttt{vietnamese\_quocngu} & L & \texttt{\string\VI} & Vietnamesisch in lateinischer Schrift \\
\texttt{vietnamese\_hannom} & L & \texttt{\string\HN} & Vietnamesische Han-Nom-Schrift \\
\bottomrule
\end{longtable}

\endgroup

\subsection{Schreibprofilkennungen (69)}
\label{app:writing-profile-ids}

Ein Schreibprofil legt die Unicode-Zeichenbereiche
und die zugehörige Satzbehandlung fest.
Die zweite Spalte zeigt seine voreingestellte
Schriftressource. Ein Strich bedeutet, dass
eine geeignete Schrift ausdrücklich gewählt
werden muss, etwa durch \texttt{[Ressource]},
\texttt{font=} oder \texttt{file=}.
Die letzte Spalte nennt die registrierten
Bereichsschlüssel. Bei mehreren Schlüsseln
wird innerhalb der Zelle umbrochen.

\begingroup\small
\begin{longtable}{@{}>{\raggedright\arraybackslash}p{0.245\linewidth}>{\raggedright\arraybackslash}p{0.305\linewidth}>{\raggedright\arraybackslash}p{0.375\linewidth}@{}}
\toprule
\textbf{Profilkennung} & \textbf{Vorgabeschrift} & \textbf{Unicode-Bereichsschlüssel} \\
\midrule
\endfirsthead
\toprule
\textbf{Profilkennung} & \textbf{Vorgabeschrift} & \textbf{Unicode-Bereichsschlüssel} \\
\midrule
\endhead
\texttt{basic} & \texttt{basic} & \texttt{basic} \\
\texttt{latin} & \texttt{basic} & \texttt{latin} \\
\texttt{greek} & \texttt{basic} & \texttt{greek-only} \\
\texttt{cyrillic} & \texttt{basic} & \texttt{cyrillic} \\
\texttt{armenian} & \texttt{armenian} & \texttt{armenian} \\
\texttt{georgian} & \texttt{georgian} & \texttt{georgian} \\
\texttt{hebrew} & \texttt{hebrew} & \texttt{hebrew} \\
\texttt{arabic} & \texttt{naskh} & \texttt{arabic} \\
\texttt{syriac} & \texttt{syriac} & \texttt{syriac} \\
\texttt{samaritan} & \texttt{samaritan} & \texttt{samaritan} \\
\texttt{aramaic} & \texttt{aramaic} & \texttt{imperial-aramaic} \\
\texttt{nabataean} & \texttt{nabataean} & \texttt{nabataean} \\
\texttt{phoenician} & \texttt{phoenician} & \texttt{phoenician} \\
\texttt{avestan} & \texttt{avestan} & \texttt{avestan} \\
\texttt{pahlavi-parthian} & \texttt{pahlavi-parthian} & \texttt{inscriptional-parthian} \\
\texttt{pahlavi-inscriptional} & \texttt{pahlavi-inscriptional} & \texttt{inscriptional-pahlavi} \\
\texttt{pahlavi-psalter} & \texttt{pahlavi-psalter} & \texttt{psalter-pahlavi} \\
\texttt{manichaean} & \texttt{manichaean} & \texttt{manichaean} \\
\texttt{sogdian} & \texttt{sogdian} & \texttt{sogdian} \\
\texttt{old-sogdian} & \texttt{old-sogdian} & \texttt{old-sogdian} \\
\texttt{old-uyghur} & \texttt{old-uyghur} & \texttt{old-uyghur} \\
\texttt{devanagari} & \texttt{devanagari} & \texttt{devanagari} \\
\texttt{tamil} & \texttt{tamil} & \texttt{tamil} \\
\texttt{thai} & \texttt{thai} & \texttt{thai} \\
\texttt{brahmi} & \texttt{brahmi} & \texttt{brahmi} \\
\texttt{kharoshthi} & \texttt{kharoshthi} & \texttt{kharoshthi} \\
\texttt{tibetan} & \texttt{tibetan} & \texttt{tibetan} \\
\texttt{coptic} & \texttt{coptic} & \shortstack[l]{\texttt{coptic-main},\\\texttt{coptic-legacy}} \\
\texttt{glagolitic} & \texttt{glagolitic} & \texttt{glagolitic} \\
\texttt{old-italic} & \texttt{old-italic} & \texttt{old-italic} \\
\texttt{old-hungarian} & \texttt{old-hungarian} & \texttt{old-hungarian} \\
\texttt{runic} & \texttt{runic} & \texttt{runic} \\
\texttt{old-turkic} & \texttt{old-turkic} & \texttt{old-turkic} \\
\texttt{carian} & \texttt{carian} & \texttt{carian} \\
\texttt{cuneiform} & \texttt{cuneiform} & \texttt{cuneiform} \\
\texttt{egyptian} & \texttt{egyptian} & \texttt{egyptian} \\
\texttt{khitan-small} & \texttt{khitan-small} & \texttt{khitan-small} \\
\texttt{tangut} & \texttt{tangut} & \texttt{tangut} \\
\texttt{han} & \texttt{shanggu} & \texttt{han} \\
\texttt{chinese} & \texttt{shanggu} & \texttt{chinese} \\
\texttt{japanese} & \texttt{japanese} & \texttt{japanese} \\
\texttt{korean} & \texttt{korean} & \texttt{korean} \\
\texttt{kana} & \texttt{japanese} & \texttt{kana} \\
\texttt{hangul} & \texttt{korean} & \texttt{hangul} \\
\texttt{bopomofo} & \texttt{iming} & \texttt{bopomofo} \\
\texttt{mongolian} & \texttt{mnglwhite} & \texttt{mongolian} \\
\texttt{jurchen} & --- & \texttt{jurchen} \\
\texttt{seal} & --- & \texttt{seal} \\
\texttt{proto-cuneiform} & --- & \texttt{proto-cuneiform} \\
\texttt{bengali} & --- & \shortstack[l]{\texttt{Bengali},\\\texttt{BengaliSupplement}} \\
\texttt{gurmukhi} & --- & \texttt{Gurmukhi} \\
\texttt{gujarati} & --- & \texttt{Gujarati} \\
\texttt{oriya} & --- & \texttt{Oriya} \\
\texttt{telugu} & --- & \texttt{Telugu} \\
\texttt{kannada} & --- & \texttt{Kannada} \\
\texttt{malayalam} & --- & \texttt{Malayalam} \\
\texttt{sinhala} & --- & \texttt{Sinhala} \\
\texttt{lao} & --- & \texttt{Lao} \\
\texttt{myanmar} & --- & \shortstack[l]{\texttt{Myanmar},\\\texttt{MyanmarExtendedA},\\\texttt{MyanmarExtendedB}} \\
\texttt{ethiopic} & --- & \shortstack[l]{\texttt{Ethiopic},\\\texttt{EthiopicExtended},\\\texttt{EthiopicExtendedA},\\\texttt{EthiopicExtendedB}} \\
\texttt{khmer} & --- & \shortstack[l]{\texttt{Khmer},\\\texttt{KhmerSymbols}} \\
\texttt{ogham} & --- & \texttt{Ogham} \\
\texttt{cherokee} & --- & \shortstack[l]{\texttt{Cherokee},\\\texttt{CherokeeSupplement}} \\
\texttt{nko} & --- & \texttt{NKo} \\
\texttt{tifinagh} & --- & \texttt{Tifinagh} \\
\texttt{yi} & --- & \shortstack[l]{\texttt{YiSyllables},\\\texttt{YiRadicals}} \\
\texttt{adlam} & --- & \texttt{Adlam} \\
\texttt{thaana} & --- & \texttt{Thaana} \\
\texttt{mandaic} & --- & \texttt{Mandaic} \\
\bottomrule
\end{longtable}

\endgroup

Die Bezeichnungen \texttt{kharoshthi} für
Schreibprofil und Ressource sowie
\texttt{kharosthi} für die frühere
Voreinstellung sind beide gültig,
gehören aber zu unterschiedlichen Namensräumen.

\subsection{Schriftressourcen und Pakete (85)}
\label{app:writing-resource-ids}

Die Kennungen dienen innerhalb von
\verb|\UseWriting| als auswählbare
Schriftressourcen. Die zweite Spalte nennt
die zugehörige Katalogfamilie. Mehrere
Kennungen können auf dieselbe Familie
und verschiedene ihrer Schnitte verweisen.

\begingroup\small
\begin{longtable}{@{}p{0.43\linewidth}p{0.47\linewidth}@{}}
\toprule
\textbf{Ressourcenkennung} & \textbf{Quellfamilie} \\
\midrule
\endfirsthead
\toprule
\textbf{Ressourcenkennung} & \textbf{Quellfamilie} \\
\midrule
\endhead
\texttt{basic} & \texttt{noto} \\
\texttt{cmu} & \texttt{cmu} \\
\texttt{times} & \texttt{times} \\
\texttt{gentium} & \texttt{gentium} \\
\texttt{charis} & \texttt{charis} \\
\texttt{libertinus} & \texttt{libertinus} \\
\texttt{latinmodern} & \texttt{mlmodern} \\
\texttt{chinese-simplified} & \texttt{chinese\_simplified} \\
\texttt{chinese-traditional} & \texttt{chinese\_traditional} \\
\texttt{japanese} & \texttt{japanese} \\
\texttt{korean} & \texttt{korean} \\
\texttt{wenjin} & \texttt{wenjin} \\
\texttt{shanggu} & \texttt{shanggu} \\
\texttt{armenian} & \texttt{armenian} \\
\texttt{devanagari} & \texttt{devanagari} \\
\texttt{tamil} & \texttt{tamil} \\
\texttt{georgian} & \texttt{georgian} \\
\texttt{tibetan} & \texttt{tibetan} \\
\texttt{hebrew} & \texttt{hebrew} \\
\texttt{arabic} & \texttt{arabic} \\
\texttt{syriac} & \texttt{syriac} \\
\texttt{syriac-eastern} & \texttt{syriac\_eastern} \\
\texttt{syriac-western} & \texttt{syriac} \\
\texttt{thai} & \texttt{thai} \\
\texttt{thai-looped} & \texttt{thai} \\
\texttt{naskh} & \texttt{arabic} \\
\texttt{kufi} & \texttt{arabic} \\
\texttt{nastaliq} & \texttt{urdu} \\
\texttt{aref} & \texttt{arabic} \\
\texttt{davidlibre} & \texttt{hebrew} \\
\texttt{khawa} & \texttt{tibetan} \\
\texttt{carian} & \texttt{anatolian} \\
\texttt{coptic} & \texttt{coptic} \\
\texttt{cuneiform} & \texttt{cuneiform} \\
\texttt{glagolitic} & \texttt{glagolitic} \\
\texttt{old-italic} & \texttt{italic} \\
\texttt{old-hungarian} & \texttt{hungarian} \\
\texttt{runic} & \texttt{runic} \\
\texttt{brahmi} & \texttt{brahmi} \\
\texttt{kharoshthi} & \texttt{kharosthi} \\
\texttt{aramaic} & \texttt{aramaic} \\
\texttt{nabataean} & \texttt{nabataean} \\
\texttt{phoenician} & \texttt{phoenician} \\
\texttt{avestan} & \texttt{avestan} \\
\texttt{pahlavi-parthian} & \texttt{pahlavi\_parthian} \\
\texttt{pahlavi-inscriptional} & \texttt{pahlavi\_inscriptional} \\
\texttt{pahlavi-psalter} & \texttt{pahlavi\_psalter} \\
\texttt{manichaean} & \texttt{manichaean} \\
\texttt{samaritan} & \texttt{samaritan} \\
\texttt{sogdian} & \texttt{sogdian} \\
\texttt{old-sogdian} & \texttt{sogdian\_old} \\
\texttt{old-turkic} & \texttt{turkic} \\
\texttt{old-uyghur} & \texttt{uyghur} \\
\texttt{egyptian} & \texttt{egyptian} \\
\texttt{khitan-small} & \texttt{khitan\_small} \\
\texttt{tangut} & \texttt{tangut} \\
\texttt{mongolian} & \texttt{mongolian} \\
\texttt{mnglwhite} & \texttt{mongolian} \\
\texttt{mngltitle} & \texttt{mongolian} \\
\texttt{mnglwriting} & \texttt{mongolian} \\
\texttt{mnglart} & \texttt{mongolian} \\
\texttt{monbaiti} & \texttt{mongolian\_baiti} \\
\texttt{abxy} & \texttt{manchu} \\
\texttt{segoeuihistoric} & \texttt{segoe} \\
\texttt{tiro} & \texttt{sanskrit} \\
\texttt{iming} & \texttt{bopomofo} \\
\texttt{zhuque} & \texttt{shanggu} \\
\texttt{simsun} & \texttt{sim} \\
\texttt{simhei} & \texttt{sim} \\
\texttt{simfang} & \texttt{sim} \\
\texttt{simkai} & \texttt{sim} \\
\texttt{arial} & \texttt{times} \\
\texttt{consolas} & \texttt{times} \\
\texttt{nomnatong} & \texttt{vietnamese\_hannom} \\
\texttt{minhnguyen} & \texttt{vietnamese\_hannom} \\
\texttt{gothicnguyen} & \texttt{vietnamese\_hannom} \\
\texttt{syrcom} & \texttt{syriac} \\
\texttt{syrcom-edessa} & \texttt{syriac} \\
\texttt{syrcom-midyat} & \texttt{syriac} \\
\texttt{syrcom-jerusalem} & \texttt{syriac} \\
\texttt{syrcom-jerusalem-bold} & \texttt{syriac} \\
\texttt{syrcom-adiabene} & \texttt{syriac\_eastern} \\
\texttt{syrcom-ctesiphon} & \texttt{syriac\_eastern} \\
\texttt{tangut-n4694} & \texttt{tangut} \\
\texttt{new-tangut-std} & \texttt{tangut} \\
\bottomrule
\end{longtable}

\endgroup

Die Tabelle beschreibt verfügbare
\emph{Ressourcennamen}, nicht den Lieferumfang
der heruntergeladenen Schriftdateien.
Je nach Installation werden vorhandene
Systemschriften oder eigens bereitgestellte
Schriftdateien verwendet. \texttt{basic}
bezeichnet die westliche Noto-Ressource;
\verb|\UseFont{noto}| bleibt die
historische Schriftvoreinstellung.

\subsection{Kurzkennungen einzelner Schriftschnitte (6)}

Diese sechs Namen sind zugleich Teil der
85 Ressourcenkennungen und können als Werte
von Schriftstiloptionen wie
\texttt{italic=} verwendet werden.
Sie bezeichnen bestimmte Schriftdateien.

\begin{longtable}{@{}p{0.26\linewidth}p{0.66\linewidth}@{}}
\toprule
\textbf{Stilkennung} & \textbf{Schriftdatei} \\
\midrule
\endfirsthead
\toprule
\textbf{Stilkennung} & \textbf{Schriftdatei} \\
\midrule
\endhead
\texttt{zhuque} & \texttt{ZhuqueFangsong-Regular.ttf} \\
\texttt{simkai} & \texttt{simkai.ttf} \\
\texttt{simfang} & \texttt{simfang.ttf} \\
\texttt{aref} & \texttt{ArefRuqaa-Regular.ttf} \\
\texttt{naskh} & \texttt{NotoNaskhArabic-Regular.ttf} \\
\texttt{nastaliq} & \texttt{NotoNastaliqUrdu-Regular.ttf} \\
\bottomrule
\end{longtable}

Beispielsweise wählt
\verb|\UseWriting{[shanggu,italic=zhuque]chinese}|
eine ausdrücklich angegebene kursive
Ersatzschrift für den entsprechenden Schnitt.

\subsection{Weitere lokale Satzbefehle}

Zusätzlich zu den Kurzbefehlen der
Familientabelle werden für senkrecht
gesetzte Schriften folgende Befehle bereitgestellt:

\begin{longtable}{@{}p{0.40\linewidth}p{0.51\linewidth}@{}}
\toprule
\textbf{Befehl} & \textbf{Darstellung} \\
\midrule
\endfirsthead
\toprule
\textbf{Befehl} & \textbf{Darstellung} \\
\midrule
\endhead
\texttt{\string\MOv[offset]\{text\}} & Mongolisch, lokal gedreht. \\
\texttt{\string\MObv[offset]\{text\}} & Mongolisch Baiti, lokal gedreht. \\
\texttt{\string\MCv[offset]\{text\}} & Mandschurisch, lokal gedreht. \\
\texttt{\string\UYv[offset]\{text\}} & Altuigurisch, lokal gedreht. \\
\bottomrule
\end{longtable}

Der optionale Versatz ist zunächst
\texttt{0pt}. Weitere Satzbefehle sind
\verb|\EG| beziehungsweise
\verb|\IMPEEgyptianText| für Ägyptisch
sowie \verb|\KHS|, \verb|\KHSstack|
und \verb|\KHSstackblock| für Kitan.

\subsection{Kompatible Klassen, Satzspiegel und Funktionen}

Die folgenden Einträge stehen für
bereits bestehende Quelldokumente bereit:

\begin{longtable}{@{}p{0.44\linewidth}p{0.46\linewidth}@{}}
\toprule
\textbf{Bisheriger Name} & \textbf{Kanonische Form} \\
\midrule
\endfirsthead
\toprule
\textbf{Bisheriger Name} & \textbf{Kanonische Form} \\
\midrule
\endhead
\texttt{impeart\_zh} & \texttt{impeart} mit \texttt{lang=zh} \\
\texttt{impereport\_zh} & \texttt{impereport} mit \texttt{lang=zh} \\
\texttt{impebook\_zh} & \texttt{impebook} mit \texttt{lang=zh} \\
\texttt{impebeamer\_zh} & \texttt{impebeamer} mit \texttt{lang=zh} \\
\texttt{nextart}, \texttt{nextreport}, \texttt{nextbook}, \texttt{nextbeamer} & Entsprechende \texttt{impe*}-Klasse \\
\texttt{next*\_zh} & Entsprechende Klasse mit \texttt{lang=zh} \\
\texttt{nextsystem.sty} & \texttt{impe.sty} \\
\bottomrule
\end{longtable}

Die früheren Layoutkennungen \texttt{report}
und \texttt{legacy\_*}, der Befehl
\verb|\UseLegacyLayout| und die
Funktionskennungen \texttt{bib}, \texttt{header}
und \texttt{lists\_envs} bleiben verfügbar.

\subsection{Kompatible Hilfsbefehle für Zusatzfunktionen}

Die bisherigen Befehle und Umgebungen
können in älteren Dokumenten weiter benutzt werden.
Für neue Dokumente stehen die regulären
\LaTeX{}-Paketbefehle im Vordergrund:

\begin{longtable}{@{}>{\raggedright\arraybackslash}p{0.43\linewidth}>{\raggedright\arraybackslash}p{0.50\linewidth}@{}}
\toprule
\textbf{Bisheriger Befehl} & \textbf{Aufgabe} \\
\midrule
\endfirsthead
\toprule
\textbf{Bisheriger Befehl} & \textbf{Aufgabe} \\
\midrule
\endhead
\shortstack[l]{\texttt{\string\Term[options]\{display\}}\\\texttt{[description]}} & Erstmaliges Hervorheben und Eintragen eines Begriffs. \\
\texttt{\string\IndexTitle} & Bisherige Registerüberschrift einstellen. \\
\texttt{TableInlineFit}, \texttt{TableLong} & Ältere Tabellenumgebungen. \\
\shortstack[l]{\texttt{TableBook*},\\\texttt{NiceBooktable*}} & Tabellengestaltung mit Anmerkungen. \\
\shortstack[l]{\texttt{OneImage*}, \texttt{PanelFigure*},\\\texttt{\string\Panel}} & Abbildungs- und Mehrfeldhilfen. \\
\texttt{ExampleBlock} & Früherer Beispielblock. \\
\bottomrule
\end{longtable}

Weitere Kompatibilitätsdetails enthält
\texttt{docs/FEATURES.md}.

\subsection{Erweiterungsschnittstellen}

Für die Registrierung eigener Schriften,
Satzspiegel und Funktionen stehen
\verb|\FontRegisterFamily|,
\verb|\LayoutPresetRegister| und
\verb|\LayoutPresetDeclare| zur Verfügung.
Weitere Angaben bieten
\texttt{docs/FONTS.md}, \texttt{docs/LAYOUTS.md},
\texttt{docs/SYSTEM.md} und \texttt{docs/EXTENDING.md}.

\section{Showcase}\label{app:showcase}

Die folgenden Seiten geben den vollständigen kanonischen IMPE-Showcase wieder.
Der Showcase ist ein visueller Begleiter des Referenzhandbuchs,
der mehrsprachiges Font-Zuordnung, komplexes Schriftformung, rechts-nach-links laufende
Schriften, vertikales Schreiben, regionale CJK-Formen und Kombinationen von
Dokumentfeatures demonstriert.

Der neueste Showcase kann über den
\href{\IMPEShowcaseURL}{aktuellen main-Branch} bezogen werden; die zu diesem
Handbuch gehörige Ausgabe wird durch den
\href{\IMPETaggedShowcaseURL}{Tag v\IMPEVersion{}} bewahrt.

\IfFileExists{\IMPEShowcaseFile}
  {\includepdf[pages=1-3,pagecommand={}]{\IMPEShowcaseFile}%
   \includepdf[pages=4-6,landscape=true,pagecommand={}]{\IMPEShowcaseFile}%
   \includepdf[pages=7-,pagecommand={}]{\IMPEShowcaseFile}}
  {\PackageError{impe-manual}
     {Erforderliche Repository-Ressource ../showcase/impe-showcase.pdf fehlt}
     {Aus manual/de kompilieren oder scripts/build_manual.ps1 aus dem Repository-Stammverzeichnis ausfuehren.}}

\end{document}
